% Property Prediction Section for Paper

\subsection{Property Prediction from H-Net Embeddings}

We evaluate H-Net embeddings as molecular and polymer property predictors, comparing against traditional RDKit descriptors (Morgan fingerprints + physicochemical properties). Features are extracted via mean pooling over token representations, then fed to XGBoost classifiers/regressors with 5-fold cross-validation.

\subsubsection{Classification Tasks}

\textbf{Blood-Brain Barrier Penetration (BBBP).} H-Net achieves AUC 0.950, outperforming RDKit (0.927) by 2.5\% absolute. This demonstrates that H-Net's learned representations capture structure-property relationships relevant to membrane permeability.

\textbf{Extended MoleculeNet Benchmarks.} We establish RDKit baselines on five additional MoleculeNet tasks:
\begin{itemize}
    \item HIV replication inhibition: AUC 0.759
    \item BACE inhibitors: AUC 0.897
    \item Clinical toxicity (ClinTox): AUC 0.865
    \item Tox21 panel (NR-AR): AUC 0.804
\end{itemize}

\subsubsection{Regression Tasks}

\textbf{Lipophilicity.} RDKit descriptors achieve MAE 0.494, outperforming H-Net (0.682) on this task. This suggests that RDKit's physicochemical descriptors provide more direct signal for lipophilicity prediction.

\textbf{Aqueous Solubility (ESOL).} RDKit baseline achieves RMSE 0.660, competitive with state-of-the-art molecular property predictors.

\textbf{Hydration Free Energy (FreeSolv).} RDKit achieves RMSE 1.131 on this challenging 642-molecule dataset.

\textbf{Polymer Properties.} For glass transition temperature (Tg), RDKit achieves MAE 24.83K while H-Net achieves 26.61K---similar performance on this challenging polymer property prediction task.

\subsubsection{Statistical Significance}

All results include 95\% bootstrap confidence intervals (Table~\ref{tab:property}). On BBBP, H-Net's improvement over RDKit is statistically significant (non-overlapping CIs).

\begin{table}[h]
\centering
\caption{Property Prediction Results Summary}
\label{tab:property}
\begin{tabular}{lllcc}
\toprule
Dataset & Domain & Metric & RDKit & H-Net \\
\midrule
\multicolumn{5}{l}{\textit{Classification (AUC-ROC)}} \\
BBBP & Molecular & AUC & 0.927 & \textbf{0.950} \\
HIV & Molecular & AUC & 0.759 & -- \\
BACE & Molecular & AUC & 0.897 & -- \\
ClinTox & Molecular & AUC & 0.865 & -- \\
Tox21 & Molecular & AUC & 0.804 & -- \\
\midrule
\multicolumn{5}{l}{\textit{Regression (lower is better)}} \\
Lipophilicity & Molecular & MAE & \textbf{0.494} & 0.682 \\
ESOL & Molecular & RMSE & 0.660 & -- \\
FreeSolv & Molecular & RMSE & 1.131 & -- \\
Tg & Polymer & MAE (K) & \textbf{24.83} & 26.61 \\
\bottomrule
\end{tabular}
\end{table}

\subsubsection{Discussion}

H-Net embeddings show promise as molecular featurizers, particularly for classification tasks where they outperform RDKit descriptors on BBBP. For regression tasks requiring precise numerical predictions, traditional descriptor-based approaches currently have an advantage.

The extended MoleculeNet baselines provide reference points for future H-Net evaluations. Running H-Net feature extraction on these datasets requires GPU compute with specialized dependencies (flash\_attn, mamba-ssm).

Key findings:
\begin{enumerate}
    \item H-Net embeddings capture meaningful structure-property relationships
    \item Mean pooling outperforms CLS pooling for aggregating token representations
    \item RDKit descriptors remain competitive baselines, especially for regression
    \item H-Net's byte-level representations generalize across molecular and polymer domains
\end{enumerate}







